\documentclass[12pt]{article}
\usepackage{amsmath, amssymb, amsthm}
\usepackage{geometry}
\geometry{margin=1in}
\usepackage{bm}

\newtheorem{definition}{Definition}
\newtheorem{lemma}{Lemma}
\newtheorem{proposition}{Proposition}
\newtheorem{remark}{Remark}

\title{The Uncertainty Component of Volatility in Veronesi (1999):\\
       Why Infrequent Regime Changes Increase Volatility}
\author{A Self-Contained Derivation}
\date{}

\begin{document}

\maketitle

\section{Introduction}

This note provides a self-contained derivation of a key result from Veronesi's (1999) paper ``Stock Market Overreaction to Bad News in Good Times: A Rational Expectations Equilibrium Model.'' Specifically, we derive the result that the ``uncertainty component'' of stock price volatility is larger when regime changes are infrequent.

The economic intuition is straightforward: when the underlying growth regime persists for long periods, investors place more weight on their current beliefs about which regime is in place. Consequently, new information causes larger revisions to expected future dividends and hence larger price movements.

\section{The Model Setup}

\subsection{The Dividend Process}

The dividend process of the risky asset is given by the stochastic differential equation:

\begin{equation}
dD_t = \theta_t dt + \sigma d\xi_t \label{eq:dividend}
\end{equation}

where:
\begin{itemize}
    \item $D_t$ is the dividend level at time $t$
    \item $\theta_t$ is an unobservable state variable representing the instantaneous dividend growth rate
    \item $\sigma > 0$ is the constant volatility of dividend innovations
    \item $\xi_t$ is a standard Wiener process
\end{itemize}

\subsection{The Hidden Regime Process}

The state variable $\theta_t$ follows a two-state, continuous-time Markov process. Let:

\begin{equation}
\theta_t \in \{\bar{\theta}, \underline{\theta}\}, \quad \text{with } \bar{\theta} > \underline{\theta}
\end{equation}

The transition probability matrix between time $t$ and $t+\Delta$ is:

\begin{equation}
\mathbf{P}(\Delta) = \begin{pmatrix}
1 - \lambda\Delta & \lambda\Delta \\
\mu\Delta & 1 - \mu\Delta
\end{pmatrix}
\end{equation}

where:
\begin{itemize}
    \item $\lambda > 0$ is the intensity (rate) of switching from the high state $\bar{\theta}$ to the low state $\underline{\theta}$
    \item $\mu > 0$ is the intensity of switching from the low state $\underline{\theta}$ to the high state $\bar{\theta}$
\end{itemize}

Thus, during an infinitesimal interval $dt$, the probability of a regime shift is $\lambda dt$ (from high to low) or $\mu dt$ (from low to high).

\subsection{Investor Beliefs}

Since investors cannot observe $\theta_t$ directly, they must infer it from the observed dividend path. Define the filtration $\{\mathcal{F}_t\}$ generated by the dividend stream $\{D_\tau\}_{\tau=0}^t$. The posterior probability of being in the high state is:

\begin{equation}
\pi_t \equiv \Pr(\theta_t = \bar{\theta} \mid \mathcal{F}_t)
\end{equation}

\subsection{The Evolution of Beliefs}

A key result from filtering theory (Liptser and Shiryayev, 1977) gives the dynamics of $\pi_t$:

\begin{lemma}[Belief Dynamics]
The posterior probability $\pi_t$ evolves according to:
\begin{align}
d\pi_t &= (\lambda + \mu)(\pi^* - \pi_t)dt + h(\pi_t)dv_t \label{eq:belief}\\
h(\pi) &\equiv \left(\frac{\bar{\theta} - \underline{\theta}}{\sigma}\right) \pi(1-\pi) \label{eq:h}\\
\pi^* &\equiv \frac{\mu}{\mu + \lambda} \label{eq:pi_star}
\end{align}
where $v_t$ is a Wiener process with respect to $\mathcal{F}_t$, defined by:
\begin{equation}
dv_t = \frac{1}{\sigma}\left[dD_t - E[dD_t \mid \mathcal{F}_t]\right]
\end{equation}
\end{lemma}

\subsection{Properties of the Belief Process}

The function $h(\pi)$ has several important properties:
\begin{itemize}
    \item $h(\pi) > 0$ for $\pi \in (0,1)$
    \item $h(\pi)$ is symmetric about $\pi = 0.5$
    \item $h(\pi)$ attains its maximum at $\pi = 0.5$, where:
    \begin{equation}
    h_{\max} = \frac{\bar{\theta} - \underline{\theta}}{4\sigma}
    \end{equation}
    \item $h(\pi) \to 0$ as $\pi \to 0$ or $\pi \to 1$
\end{itemize}

This implies that investor uncertainty about the true state is maximized when $\pi = 0.5$ (maximum uncertainty about which regime prevails) and minimized when $\pi$ is near 0 or 1 (investors are nearly certain).

\section{The Risk-Neutral Price}

\subsection{Definition}

The risk-neutral price $P^*(D_t, \pi_t)$ is defined as the discounted expected value of all future dividends, evaluated at the risk-free rate $r$:

\begin{equation}
P^*(D_t, \pi_t) \equiv E_t\left[\int_0^\infty e^{-rs} D_{t+s} ds\right]
\end{equation}

This is the price that would prevail in a risk-neutral world where investors do not require a risk premium.

\subsection{Closed-Form Expression}

\begin{lemma}[Risk-Neutral Price]
The risk-neutral price is given by:
\begin{align}
P^*(D, \pi) &= \frac{1}{r}D + \frac{1}{r^2}\theta^* + \frac{\bar{\theta} - \underline{\theta}}{r(\lambda + \mu + r)}(\pi - \pi^*) \label{eq:risk_neutral}
\end{align}
where $\theta^* = \pi^*\bar{\theta} + (1-\pi^*)\underline{\theta}$ is the unconditional expected growth rate.
\end{lemma}

\begin{proof}
First, compute $E_t[\theta_{t+u}]$. Using the transition probabilities of the Markov chain:
\begin{equation}
E_t[\theta_{t+u}] = \theta^* - (\bar{\theta} - \underline{\theta})(\pi^* - \pi_t)e^{-(\lambda+\mu)u}
\end{equation}

Then:
\begin{align}
E_t[D_{t+s}] &= D_t + \int_0^s E_t[\theta_{t+u}]du \\
&= D_t + \theta^* s + (\bar{\theta} - \underline{\theta})(\pi^* - \pi_t)\left(\frac{e^{-(\lambda+\mu)s} - 1}{\lambda + \mu}\right)
\end{align}

Finally:
\begin{align}
P^* &= \int_0^\infty e^{-rs} E_t[D_{t+s}]ds \\
&= \frac{D_t}{r} + \frac{\theta^*}{r^2} - \frac{\bar{\theta} - \underline{\theta}}{r(\lambda + \mu + r)}(\pi^* - \pi_t)
\end{align}
\end{proof}

The key observation is that $P^*$ is linear in $\pi$ with slope:
\begin{equation}
p_\pi \equiv \frac{\partial P^*}{\partial \pi} = \frac{\bar{\theta} - \underline{\theta}}{r(\lambda + \mu + r)} \label{eq:p_pi}
\end{equation}

\section{The Uncertainty Component of Volatility}

\subsection{Definition}

The volatility of the risk-neutral price, which Veronesi calls the ``uncertainty component,'' is defined as:
\begin{equation}
\sigma_{P^*}^2(\pi) \equiv \frac{E_t[(dP^*)^2]}{dt}
\end{equation}

\subsection{Derivation}

Apply Itô's lemma to $P^*(D, \pi)$. Since $P^*$ depends on both $D$ and $\pi$:
\begin{equation}
dP^* = \frac{\partial P^*}{\partial D}dD + \frac{\partial P^*}{\partial \pi}d\pi + \frac{1}{2}\frac{\partial^2 P^*}{\partial \pi^2}(d\pi)^2
\end{equation}

But from equation \eqref{eq:risk_neutral}, $\frac{\partial P^*}{\partial D} = \frac{1}{r}$ and $\frac{\partial P^*}{\partial \pi} = p_\pi$, while $\frac{\partial^2 P^*}{\partial \pi^2} = 0$. Therefore:
\begin{equation}
dP^* = \frac{1}{r}dD + p_\pi d\pi
\end{equation}

Substituting $dD$ from \eqref{eq:dividend} and $d\pi$ from \eqref{eq:belief}:
\begin{align}
dP^* &= \frac{1}{r}(\theta_t dt + \sigma dv_t) + p_\pi\left[(\lambda+\mu)(\pi^*-\pi)dt + h(\pi)dv_t\right] \\
&= \left[\frac{1}{r}\theta_t + p_\pi(\lambda+\mu)(\pi^*-\pi)\right]dt + \left[\frac{\sigma}{r} + p_\pi h(\pi)\right]dv_t
\end{align}

The instantaneous volatility is the coefficient on $dv_t$:
\begin{equation}
\sigma_{P^*}(\pi) = \frac{\sigma}{r} + p_\pi h(\pi) \label{eq:vol_uncertainty}
\end{equation}

Substituting $p_\pi$ from \eqref{eq:p_pi} and $h(\pi)$ from \eqref{eq:h}:
\begin{equation}
\boxed{\sigma_{P^*}(\pi) = \frac{\sigma}{r} + \frac{(\bar{\theta} - \underline{\theta})^2}{r(r + \lambda + \mu)\sigma} \pi(1-\pi)} \label{eq:vol_final}
\end{equation}

\section{Comparative Statics: Why Infrequent Changes Increase Volatility}

\subsection{The Effect of $\lambda + \mu$}

From equation \eqref{eq:vol_final}, the coefficient on $\pi(1-\pi)$ is:
\begin{equation}
C \equiv \frac{(\bar{\theta} - \underline{\theta})^2}{r(r + \lambda + \mu)\sigma}
\end{equation}

Since $C$ is strictly decreasing in $(\lambda + \mu)$, we have:
\begin{equation}
\frac{\partial \sigma_{P^*}(\pi)}{\partial (\lambda + \mu)} = -\frac{(\bar{\theta} - \underline{\theta})^2}{r(r + \lambda + \mu)^2\sigma} \pi(1-\pi) < 0
\end{equation}

Thus, for any $\pi \in (0,1)$, a smaller $(\lambda + \mu)$ (i.e., infrequent regime changes) implies a larger $\sigma_{P^*}(\pi)$.

\begin{proposition}[Effect of Infrequent Regime Changes]
When regime changes are infrequent ($\lambda + \mu$ is small), the uncertainty component of volatility is larger, ceteris paribus.
\end{proposition}

\subsection{Intuition}

The parameter $p_\pi = \frac{\bar{\theta} - \underline{\theta}}{r(\lambda + \mu + r)}$ measures the sensitivity of the risk-neutral price to changes in beliefs. When $\lambda + \mu$ is small:
\begin{itemize}
    \item The hidden state $\theta_t$ is highly persistent
    \item Current beliefs are more informative about the distant future
    \item A given belief update $\Delta\pi$ causes a larger revision in expected future dividends
    \item Therefore, price volatility is higher
\end{itemize}

In the limiting case as $\lambda + \mu \to 0$ (no regime changes ever occur), we have:
\begin{equation}
p_\pi \to \frac{\bar{\theta} - \underline{\theta}}{r^2}
\end{equation}
which is the maximum possible sensitivity.

\subsection{Other Comparative Statics}

From the same expression \eqref{eq:vol_final}, we can derive two additional predictions:

\subsubsection{Magnitude of Regime Differences}
\begin{equation}
\frac{\partial \sigma_{P^*}(\pi)}{\partial (\bar{\theta} - \underline{\theta})} = \frac{2(\bar{\theta} - \underline{\theta})}{r(r + \lambda + \mu)\sigma} \pi(1-\pi) > 0
\end{equation}
More dramatic differences between the two regimes increase volatility.

\subsubsection{Risk-Free Rate}
\begin{equation}
\frac{\partial \sigma_{P^*}(\pi)}{\partial r} = -\frac{\sigma}{r^2} - \frac{(\bar{\theta} - \underline{\theta})^2 (2r + \lambda + \mu)}{r^2(r + \lambda + \mu)^2\sigma} \pi(1-\pi) < 0
\end{equation}
A lower discount rate $r$ increases volatility, because future cash flows are more heavily weighted.

\section{Connection to the Full Model}

In the complete model with risk-averse investors, total price-change volatility includes an additional term:
\begin{equation}
\sigma_P(\pi) = \sigma_{P^*}(\pi) + S'(\pi)h(\pi)
\end{equation}
where $S(\pi)$ is a negative, convex function arising from the hedging demands of risk-averse investors. This risk-aversion component introduces asymmetry: volatility is amplified in good times ($\pi > 0.5$) and dampened in bad times ($\pi < 0.5$). However, the uncertainty component $\sigma_{P^*}(\pi)$ remains the foundation and exhibits the comparative statics derived above.

\section{Summary}

The uncertainty component of volatility is:
\begin{equation}
\sigma_{P^*}(\pi) = \frac{\sigma}{r} + \underbrace{\frac{(\bar{\theta} - \underline{\theta})^2}{r(r + \lambda + \mu)\sigma}}_{\text{amplitude}} \pi(1-\pi)
\end{equation}

The amplitude increases when:
\begin{itemize}
    \item Regime changes are infrequent ($\lambda + \mu$ small)
    \item Regime differences are large ($\bar{\theta} - \underline{\theta}$ large)
    \item The risk-free rate is low ($r$ small)
\end{itemize}

The economic mechanism is that $p_\pi$, the sensitivity of the risk-neutral price to beliefs, is larger when regimes persist longer, causing the same news to generate larger price movements.

\begin{remark}
Throughout this derivation, we have assumed a fixed supply of the risky asset (normalized to 1) and a constant risk-free rate $r$. These assumptions, standard in the literature, allow us to isolate the effect of learning and regime shifts on volatility.
\end{remark}

\end{document}